\section{Phase estimates before key collisions}
\label{sec:phases}

The estimates in this section supply the compatible small tables
used in \Cref{prop:status-preparation}.  Channel thinning makes the
retained mass and the compatibility margin positive constants chosen
before the channel size.  The original frame law is \(\rawlaw\),
whereas all marginal bounds below are against the thinned law \(\pi\).

Throughout this section, different units are independent, but the two
endpoints of one unit may have arbitrary dependence.  A marked unit is
a unit together with choices
\(\lambda_i\in C_i^{\mathrm{old}}\), \(i=1,2\), where
\(C_i^{\mathrm{old}}\) denotes its effective space after the first
peeling.  Marks are permitted to depend on the entire unit.  Put
\begin{equation}
 \Delta_O=U_1\lambda_1+U_2\lambda_2,\qquad
 u_{O,S}=\sum_{i\in S}u_i\quad(S\subseteq\{1,2\}).
 \label{eq:delta-phase}
\end{equation}
The total component rank of \(\Delta_O\) is at most
\(r=2K_1+2\).  We use \(L_0,d_0,u_0,K\) from
\Cref{eq:law-budgets,eq:pin-budgets}; in particular,
\[
 L_0=\lceil100(D+10)\rceil,\qquad d_0=4rL_0,\qquad
 u_0=10(K_1+d_0+1).
\]
Write \(\varepsilon_{\rm ch}=10|\cE|2^h\rho\), as in
\Cref{lem:channel-thinning}.  A marked law with endpoint marginals
at most \(\Lambda\pi\) has unchanged primal marginal bound
\(\Lambda\operatorname{pr}_*\rawlaw\), but its full-frame bound is only
\(\Lambda2^{\varepsilon_{\rm ch}N}\rawlaw\).
Here \(\Lambda\) is a fixed constant specified before \(h\).
In the coloured test it comes from the pooled law and the total mass
retained, and does not depend on the probability of an individual
colour.  The starting absolute pair cap is \(2^{3000gN}\pi^2\).
After the first peeling and restrictions of fixed positive mass it
is at most \(2^{(D+.01)N}\rawlaw^2\), by
\Cref{lem:channel-thinning} and \(D=4000g\).
Every use of a full-frame image bound below explicitly pays the
loss \(\varepsilon_{\rm ch}N\).  We choose
\(\varepsilon_{\rm ch}<\min\{\zeta/100,10^{-4}\}\), after \(h\),
and choose \(M_0\) afterward to make each primal ratio loss as small
as required.  A fixed multiplicative constant is absorbed by taking
\(n\) sufficiently large.

\begin{definition}
A componentwise cover is a collection of subspaces
\(T_e^\pm\subseteq V_e^\pm\).  A tensor \(\xi=(\xi_e)_e\) is covered
by it if
\[
 \xi_e\in T_e^+\otimes V_e^-+V_e^+\otimes T_e^-
 \quad\text{for every }e.
\]
The dimensions of the two modes of the cover are
\(\sum_e\dim T_e^+\) and \(\sum_e\dim T_e^-\).
\end{definition}

\begin{lemma}[Rank growth without a cover]
\label{lem:no-cover-rank}
Let \(r\ge1\) be an integer, set \(m_r=100r+10\) and
\(q_r=2^{-m_r}\), and let \(s\) be a power of two with
\(s\ge2^{m_r}\).  Suppose \(T_1,\ldots,T_s\) are independent copies
of a random tensor of total component rank at most \(r\).
If every fixed componentwise cover with dimension at most \(rs\)
in each mode has probability less than \(p\), then
\begin{equation}
 \PP\left\{\rank\left(\sum_{j=1}^sT_j\right)<q_rs/4\right\}
 \le 4^s p^{q_rs/2}.
 \label{eq:no-cover-rank}
\end{equation}
Here rank means the sum of component ranks.
\end{lemma}

\begin{proof}
Regard all components as blocks in the direct sums of the two mode
spaces.  Their column and row spaces are direct sums of
componentwise spaces, so that every cover obtained in this proof is
componentwise.  We first make a deterministic observation about any
realized list of tensors.

Expose a subset \(E\) of the indices and quotient the two modes by
the column and row spans of the exposed tensors.  If \(b=s-|E|\)
indices remain, let \(C_j,R_j\) be their projected original column
and row spaces.  Define their span deficits by
\[
 \delta_C=\sum_{j\notin E}\dim C_j-\dim\sum_{j\notin E}C_j,\qquad
 \delta_R=\sum_{j\notin E}\dim R_j-\dim\sum_{j\notin E}R_j.
\]
The projected sum has rank at least
\begin{equation}
 \#\{j\notin E:\overline T_j\ne0\}-\delta_C-\delta_R.
 \label{eq:rank-deficits}
\end{equation}
Indeed, before adding the mode spaces, place the projected tensors
on the separate formal summands \(C_j\otimes R_j\).  Their block
diagonal sum has rank at least the displayed number of nonzero
tensors.  Applying the two addition maps can decrease rank by at
most their kernel dimensions, which are \(\delta_C,\delta_R\).

There is an exposed set leaving \(b=s/2^j\) indices for some
\(0\le j<m_r\) such that
\(\delta_C+\delta_R\le b/4\).  To prove this, order the realized
direction spaces by a uniform random permutation.  In either mode,
let \(a_t\) be the expected increase of their span dimension at the
\(t\)-th step.  Submodularity and exchangeability give
\[
 r\ge a_1\ge a_2\ge\cdots\ge a_s\ge0.
\]
For \(b_j=s/2^j\), set
\[
 f_j=a_{s-b_j+1},\qquad
 v_j=f_j-\frac1{b_j}\sum_{t=s-b_j+1}^s a_t.
\]
Each remaining index has expected projected dimension \(f_j\)
immediately after the first \(s-b_j\) exposures.  Consequently its
mode's expected deficit divided by \(b_j\) is \(v_j\).
The average increment on the first half of this tail is at least
\(f_{j+1}\).  It follows that
\[
 v_j\le f_j-f_{j+1}+\tfrac12v_{j+1}.
\]
Summing for \(0\le j<m_r\) yields
\[
 \sum_{j=0}^{m_r-1}v_j
 \le2(f_0-f_{m_r})-v_0+v_{m_r}\le2r.
\]
The two modes therefore have total expected deficit ratios at most
\(4r\) across these \(m_r\) scales.  Some scale has expected ratio
less than \(1/4\), and some permutation realizes a ratio at most
\(1/4\), as claimed.

If the original sum has rank less than \(q_rs/4\), its projection
at this exposed set has no larger rank.  By
\Cref{eq:rank-deficits}, at most
\(q_rs/4+b/4\le b/2\) remaining tensors have nonzero projections.
At least \(q_rs/2\) of them thus belong to the cover formed by the
exposed column and row spans.  Both cover dimensions are at most
\(rs\).

For fixed disjoint index sets \(E,J\), with
\(|J|=\lceil q_rs/2\rceil\), condition only on the tensors indexed
by \(E\).  The cover is then fixed, and the tensors indexed by \(J\)
are still independent.  The probability that all are covered is
at most \(p^{|J|}\).  There are at most \(2^s\) choices for each
index set.  Taking their union proves \Cref{eq:no-cover-rank}.
\end{proof}

\begin{theorem}[Phase alternative]
\label{thm:phase-alternative}
Suppose a marked-unit law \(\sigma\) satisfies
\[
 \sigma_i\le\Lambda\pi\quad(i=1,2),\qquad
 \sigma\le2^{(D+.01)N}\rawlaw^2,
\]
where \(\Lambda\) is fixed before the channel size, and suppose
\(\Delta_O\ne0\) throughout.  For all sufficiently large \(n\),
one can retain a sublaw of relative mass at least
\(\delta_{\rm ph}>0\), depending only on the early parameters,
such that one of the following holds for independent draws \(A,B\):
\begin{enumerate}[label=\textup{(\roman*)}]
\item every nontrivial phase mean
\[
 \E(-1)^{u_{B,S}(\Delta_A)+u_{A,R}(\Delta_B)},
 \qquad(S,R)\ne(\varnothing,\varnothing),
\]
has absolute value less than \(.005\);
\item writing \(F(A,B)=u_{B,\{1,2\}}(\Delta_A)\), one has
\(F(A,B)=0\) for every pair of retained units, and every displayed
mean for which \(|S|=1\) or \(|R|=1\) is \(2^{-\Omega(N)}\).
\end{enumerate}
All marks remain in their first effective spaces.  No additional
image directions are pinned.  Neither the retained-mass bound nor
the tiny-cover dimension used in the proof depends on selector or
channel sizes; the lower bound required of \(h\) may depend on
\(\Lambda\).
\end{theorem}

\begin{proof}
We first choose the parameters and treat the case with no large cover.
In the cover case we pass to offset records. A tiny cover either makes
\(F\) identically zero or permits a gain from thinning; in the absence
of a tiny cover, a fresh-direction count gives mixing.

Write
\[
 c_s=10^{-12}(1+r)^{-4},\qquad q_0=2^{-100r-10},
 \qquad \theta=10^{-4},\qquad \Gamma=D+2r+10,
\]
and let \(s\) be the largest power of two not exceeding \(c_sN\).
For large \(n\), it satisfies the size requirement of
\Cref{lem:no-cover-rank} and \(s\ge c_sN/2\).
Choose \(p_0>0\), using only these early parameters, so small that
\begin{equation}
 \frac{q_0}{2}\log_2(1/p_0)
 \ge 2+\log_2(1/\theta)+\frac{2\Gamma}{c_s}.
 \label{eq:early-p0}
\end{equation}
When the channel size is chosen, require also
\begin{equation}
 \frac{hq_0}{4}
 \ge \log_2(1/\theta)+\frac{2\Gamma}{c_s}.
 \label{eq:early-h-moment}
\end{equation}
Further lower bounds for \(h\) arising below will again involve
only early parameters.

\medskip\noindent
\emph{No large cover.}
Suppose no fixed cover of dimension at most \(rs\) in each mode
has mass at least \(p_0\).  Fix a candidate value \(\xi\) of
\(\Delta_B\), a nonempty \(S\), and a set \(R\).  Define the
unit-side sign \(\psi(A)=(-1)^{u_{A,R}(\xi)}\).  Against completely
iid channel columns at the endpoints selected by \(S\), expansion
of the even \(s\)-th moment gives
\begin{align}
 &\E_{\mathrm{iid\ channels}}
 \left|\E_A\psi(A)(-1)^{u_{B,S}(\Delta_A)}\right|^s
 \notag\\
 &\hspace{25mm}\le
 \E_{A_1,\ldots,A_s}
       2^{-h\rank(\Delta_{A_1}+\cdots+\Delta_{A_s})}.
 \label{eq:phase-moment}
\end{align}
For one channel pair the expectation of its bilinear character
on a tensor \(\xi\) is \(2^{-\rank\xi}\).  Independence over
channels and components proves this assertion; selecting both
endpoints only squares the factor.  The signs \(\psi(A_j)\)
have absolute value one and can be removed for an upper bound.
By \Cref{lem:no-cover-rank}, the last expression is at most
\[
 2^{2s}p_0^{q_0s/2}+2^{-hq_0s/4}.
\]
The channel marginal comparison in \Cref{lem:frame-images}
multiplies this bound by a constant depending on \(h\) and the
number of components, but independent of \(n\).
Markov's inequality and
\Cref{eq:early-p0,eq:early-h-moment} bound the probability that
the inner mean has magnitude at least \(\theta\) by
\(C_h2^{-\Gamma N}\).

There are at most \(C_r2^{2rN}\) tensors of total component rank
at most \(r\): choose their component ranks, then factor each
matrix through that rank.  The component-rank choices and the
overcounting factors are bounded independently of \(n\).
Take a union over these candidate values and the bounded choices
of \(S,R\), and then pay the joint-law density
\(2^{(D+.01)N}\).  The exceptional \(B\)-mass is exponentially
small.  On all other \(B\)'s the conditional \(A\)-mean is at most
\(\theta\), even after inserting the actual value \(\Delta_B\).
Every required phase mean is therefore at most
\(\theta+2^{-\Omega(N)}<.005\).  Interchanging \(A,B\) deals
with \(S=\varnothing\).  This proves the first alternative
without a further restriction.

\medskip\noindent
\emph{A large cover and its offsets.}
Otherwise retain mass at least \(p_0\) in a fixed cover
\((T_e^+,T_e^-)\) of dimension at most \(rs\) in each mode.
Each individual primal frame span avoids these fixed subspaces
except on exponentially small mass.  Indeed its dimension is
\(\dim\cB\), its location is uniform under \(\rawlaw\), and
\[
 \PP_{\rawlaw}\{\im P_{ie}\cap T_e^+\ne0\}
 \le 2^{\dim\cB+\dim T_e^+-N+1};
\]
the same estimate holds for the minus mode.  Here
\(rs\le rc_sN\) and \(\dim\cB/N\) is sufficiently small.
The primal marginal bound \(\Lambda\operatorname{pr}_*\rawlaw\) permits
deletion of these exceptions; this step does not use a full-frame
bound against \(\rawlaw\).  Normalize the retained law, still denoted \(\sigma\).
Its marginal bound \(\sigma_i\le L\pi\) has
\(L\le2\Lambda/p_0\) for all sufficiently large \(n\).
Thus this fixed bound may also depend on the given \(\Lambda\),
whereas the retained-mass bound remains independent of \(\Lambda\).

Project both modes of each component modulo \(T_e^\pm\).
Because \(\Delta\) is covered, the two endpoint tensors have the
same quotient.  The projections are injective on the individual
primal spans, so they preserve each tensor rank.  Choose a shortest
outer-product factorization of the common quotient tensor and
lift each factor to the two endpoint primal spans.  Componentwise,
the difference has the form
\begin{equation}
 \Delta=\sum_\alpha
   \bigl(p_\alpha b_\alpha\trans+
         a_\alpha q_\alpha\trans+
         a_\alpha b_\alpha\trans\bigr),
 \quad
 a_\alpha\in T^+,\quad b_\alpha\in T^- .
 \label{eq:cover-lifting}
\end{equation}
The first endpoint is the reference here; its factors are
\(p_\alpha,q_\alpha\), and those at the second endpoint are
\(p_\alpha+a_\alpha,q_\alpha+b_\alpha\).
The factor lists at either endpoint are independent in their
respective modes.

Choose bases for the spans of the offsets \(a_\alpha,b_\alpha\)
in each component.  Their combined ordered list is denoted
\(\omega_O\), and its size is \(t_O\), with
\begin{equation}
 1\le t_O\le2r.
 \label{eq:offset-rank}
\end{equation}
The lower bound follows from \(\Delta\ne0\).
Collect the companions of these offset basis vectors in a tuple
\(Z_O\): minus companions pair with plus offsets and vice versa.
The companions are independent within each component and mode.
For example, writing \(b_\alpha=\sum_kc_{k\alpha}\omega_k^-\)
gives a full-row-rank matrix \((c_{k\alpha})\), so the vectors
\(\sum_\alpha c_{k\alpha}p_\alpha\) are independent.
Using the second endpoint as reference changes each companion
by a known combination of offsets, and preserves independence.

An offset record consists of these basis vectors, their component
and sign assignments, and the bounded coefficient arrays in
\Cref{eq:cover-lifting}.  Its logarithmic count is at most
\[
 2r^2s+O_{r,g}(1)<.002N
\]
for sufficiently large \(n\).  Choosing a nominal representation
of a companion tuple in one endpoint's primal frame costs
\(2^{O(r\dim\cB)}=2^{O(n)}\) possibilities.  This last cost has
an arbitrarily small ratio exponent after choosing \(M_0\).
Fix deterministic conventions for all these records and
factorizations; their choices can be arbitrary functions of the
marked unit.

For a linear tensor functional \(u\), let \(\Phi_u(\omega)\)
be the tuple obtained by contracting \(u\) with the offset basis
vectors.  In particular, on a plus offset \(a\) its corresponding
plus output is \(X(Y\trans a)\); on a minus offset \(b\) the
minus output is \(Y(X\trans b)\).  At fixed two offset records
the phase in the theorem is, up to a sum of two unary phases,
\begin{equation}
 Z_A\cdot\Phi_{u_{B,S}}(\omega_A)
 +Z_B\cdot\Phi_{u_{A,R}}(\omega_B).
 \label{eq:offset-walsh}
\end{equation}
The unary phases are evaluations of the last term of
\Cref{eq:cover-lifting}; fixed offset records make them separate
functions of \(A\) and \(B\).

\medskip\noindent
\emph{Singleton phase estimates.}
All record laws in the following calculation are unnormalized
subprobabilities of \(\sigma\).  For any record, the companion
tuple alone has point masses at most
\begin{equation}
 p_{\max}(Z_O)\le2^{-t_ON+.01N}.
 \label{eq:companion-pmax}
\end{equation}
To obtain this, enumerate its nominal primal representations,
apply the individual primal frame-image estimate and the unchanged
primal marginal bound \(L\operatorname{pr}_*\rawlaw\), and choose \(M_0\) large enough to absorb the
coefficient count and lost primal codimensions.

Suppose \(S=\{j\}\), and use endpoint \(j\) as the reference
for the \(B\)-companions.  For fixed \(\omega_A\), discard
pairs on which the coefficient lists
\(Y_j\trans a\) or \(X_j\trans b\), for its independent offsets,
fail full rank in any component.  Under the original law, comparison with iid columns bounds the
transpose-rank failure on \(d\) independent ambient vectors by
\(2^{d-h+1}\) for large \(n\).  The fixed-list assertion of
\Cref{lem:channel-thinning} bounds the same event under \(\pi\)
by \(2^{d-h+1}+o(1)\), uniformly over the fixed list.  For each
fixed \(A\), the list being tested is fixed before drawing \(B\).
Consequently the discarded pair mass is at most
\[
 L\,O(r)\,2^{2r-h}+o(1)<10^{-4},
\]
on making \(h\) sufficiently large in terms of the fixed marginal
cap and the early rank bound.

On the remaining pairs, the tuple
\((Z_B,\Phi_{u_{B,j}}(\omega_A))\) has point masses at most
\begin{equation}
 2^{-(t_B+t_A)N+(.02+\varepsilon_{\rm ch})N}.
 \label{eq:singleton-joint-pmax}
\end{equation}
Here is a direct conditional count.  Enumerate a nominal
representation of \(Z_B\), first test its values, and condition
on the primal frames at endpoint \(j\).  The two channel maps
are independent uniform-column draws in the appropriate
annihilators, apart from a bounded injectivity conditioning.
Freeze the values of \(Y\trans a\) and \(X\trans b\).  Their
number of possibilities is bounded in \(n\), and their lists
have full rank on the retained event.  Ignore the equations
that originally defined these coefficient values.  Prescribing
their images under \(X\) and \(Y\) then costs at least
\((N-\dim\cB)t_A\) bits, up to a bounded factor.
This is a count under \(\rawlaw\), conditional on its primal
frames.  Replacing its channel law by that of \(\pi\) costs at
most \(2^{\varepsilon_{\rm ch}N}\); applying the endpoint marginal
bound costs the fixed factor \(L\).  Together with the primal
companion test this gives \Cref{eq:singleton-joint-pmax}.  The ignored equations only
enlarge the event, so no independence between a coefficient
list and the map defining it has been assumed.

Apply \Cref{lem:walsh} to \Cref{eq:offset-walsh}, on
\((t_A+t_B)N\) bits.  On the \(A\)-side the point-mass bound for
the entire tuple is no larger than
\Cref{eq:companion-pmax}; on the \(B\)-side use
\Cref{eq:singleton-joint-pmax}.  Separate unary phases have
absolute value one.  The contribution for two fixed records is
at most
\[
 2^{(t_A+t_B)N/2}
 \left(2^{-t_AN+.01N}
       2^{-(t_A+t_B)N+(.02+\varepsilon_{\rm ch})N}\right)^{1/2}
 =2^{(-t_A/2+.015+\varepsilon_{\rm ch}/2)N}.
 \]
There are at most \(2^{.004N}\) record pairs.  By
\Cref{eq:offset-rank}, their total is at most \(2^{-.480N}\).
Thus every singleton phase mean is bounded by
\(10^{-4}+2^{-.480N}\).  The same proof applies when
\(|R|=1\).

\medskip\noindent
\emph{A tiny cover or mixing of the full endpoint sum.}
Put \(\delta=.0005\).  If a cover of total dimension at most
\(d_0\) contains the whole offset list with probability at least
\(\delta/4\), restrict to those units.  Intersect its spaces componentwise with the preceding large cover.
This preserves every captured offset list and does not increase its
dimension.  Each individual primal span therefore still avoids it.  Delete units whose channel transpose at either
endpoint is not injective on a fixed basis of any component and
sign of this tiny cover.  By \Cref{lem:channel-thinning} and the
endpoint bound, the deleted relative mass is at most
\[
 L'\,O(d_0+|\cE|)2^{d_0-h}+o(1),
 \qquad L'\le 4L/\delta.
\]
Take \(h\) so this is less than \(1/2\).  This leaves relative mass
at least \(\delta/8\) in the large-cover law.  All future offset
lists in this tiny cover now pass the singleton rank tests.

Repeat \Cref{eq:cover-lifting} with the tiny cover as the quotient
directions.  Fix an offset record: the actual offset basis, its
component and sign assignments, and its bounded binary coefficient
arrays.  These data have at most
\[
 R_0=2^{100(r+1)^2(d_0+|\cE|+1)}
\]
possibilities.  Indeed at most \(2r\) vectors are chosen from a
fixed space of dimension \(d_0\), at most \(2r\) component/sign
assignments are made, and all coefficient arrays have bounded
sizes in terms of \(r,|\cE|\).  The deliberately generous exponent
bounds the product of these finite counts.  Retain a record of
mass at least \(1/R_0\).  This fixes a common ordered tuple
\(\omega\), of length \(1\le t\le2r\), but fixes neither the
companions nor any frame image.  In particular this restriction
is not an additional pinning step.

Split once more, retaining at least half the mass, according as
\(\Phi_{u_{O,\{1,2\}}}(\omega)\) is zero or nonzero.  All
normalizations in this paragraph are bounded by early constants.
In the zero case each term in \Cref{eq:cover-lifting} has one
factor from \(\omega\), so
\[
 u_{B,\{1,2\}}(\Delta_A)=0
 \quad\hbox{for every retained } A,B.
\]
This includes the last, offset--offset term of the lift.  The
singleton count already proved, with no rank discard because the
whole tiny cover passes the transpose tests, gives exponentially
small singleton means under this fixed restriction.  Thus
alternative \textup{(ii)} holds.

In the nonzero case we claim the stronger full-sum point bound
\begin{equation}
 \PP_\sigma\{Z_O=z,
       \Phi_{u_{O,\{1,2\}}}(\omega)=\phi\}
 \le 2^{-tN-\rho N/3}
 \quad(\phi\ne0).
 \label{eq:tp-fullsum-pmax}
\end{equation}
Here \(\sigma\) denotes the normalized retained law; its endpoint
marginals are at most \(L''\pi\) for an early constant \(L''\).
To verify the claim, choose endpoint one as reference in the lift.
Enumerate all nominal primal descriptions of its \(t\) independent
companions; their count is \(2^{O(n)}\).  For each description,
prescribing the companion images \(z\) has probability at most
\(2^{-tN+O(n)}\) under the primal marginal of \(\pi\), which is
exactly the primal marginal of \(\rawlaw\).
Choose a nonzero coordinate \(d\) of the prescribed output \(\phi\).
The corresponding contraction of each endpoint's channel form is
a nonzero channel combination, because its transpose is injective
on the fixed cover basis.  Call these combinations \(v_1,v_2\).
They lie in the same prescribed component/sign subset, and
\(v_1+v_2=d\ne0\).  Necessarily, therefore, endpoint one has a
nonzero channel combination in the translate by \(d\) of its
allowed subset.  Conditional on all its primal frames, this
endpoint-one event has \(\pi\)-probability at most
\(2^{-\rho N/2}\), by \Cref{lem:channel-thinning}.

This is a necessary event on endpoint one alone.  Consequently
its probability under the dependent two-endpoint law is bounded
using only the endpoint marginal \(L''\pi\): no independence of
the endpoints, and no conditional density bound for endpoint two,
is used.  Sum the nominal descriptions to obtain
\[
 \PP_\sigma\{Z_O=z,\Phi=\phi\}
 \le L''2^{-tN+C n}2^{-\rho N/2}
\]
for a constant \(C\) fixed before \(M_0\).
Choose \(M_0\) so \(C/M_0<\rho/12\), and then \(n\) so
\(\log_2 L''/N<\rho/12\).  This proves
\Cref{eq:tp-fullsum-pmax}.

For the other side of the Walsh estimate the companion-only
bound can now be sharpened to
\begin{equation}
 p_{\max}(Z_O)\le2^{-tN+\rho N/12}.
 \label{eq:tp-fine-companion-pmax}
\end{equation}
Its errors are solely \(O(n)\) primal image and nominal-description
costs plus a fixed normalization, so increasing \(M_0\) after
\(\rho\) proves this bound.  There is no full-frame thinning loss
in \Cref{eq:tp-fine-companion-pmax}.
For a phase containing a full-sum term, apply \Cref{lem:walsh} to
\Cref{eq:offset-walsh}, using
\Cref{eq:tp-fullsum-pmax} on that term's side and
\Cref{eq:tp-fine-companion-pmax} on the other.  Both records are
already fixed, so there is no exponential record union.  The
absolute value is at most
\[
 2^{2tN/2}
 \left(2^{-tN+\rho N/12}2^{-tN-\rho N/3}\right)^{1/2}
 =2^{-\rho N/8}.
\]
When the other phase term is absent, its output coordinates are
zero; the companion bound still bounds the full tuple.  Phases
containing a singleton have the previous exponential estimate.
Every nontrivial phase is of one of these types, so the nonzero
case gives alternative \textup{(i)}.

Assume instead that no such cover captures mass \(\delta/4\).
Fix an \(A\)-offset list \(v=\omega_A\), and use endpoint \(1\)
as reference throughout the next calculation.  Call \(B\)
frequent for this list if its offset record \(m\), its companion
value \(z\), and its output \(\phi=\Phi_{u_{B,\{1,2\}}}(v)\)
have joint point probability greater than
\(2^{-(t_B+.5)N}\).  For fixed \(m,z\), define the deterministic
list
\begin{equation}
 {\cal F}_v(m,z)=
 \left\{\phi:
  \sigma\bigl(m_B=m,Z_B=z,
       \Phi_{u_{B,\{1,2\}}}(v)=\phi\bigr)
       >2^{-(t_B+.5)N}\right\}.
 \label{eq:frequent-list}
\end{equation}
By \Cref{eq:companion-pmax},
\(\abs{{\cal F}_v(m,z)}\le2^{.51N}\).

We claim that the frequent-pair mass is at most \(\delta\).
Otherwise a set of \(B\)'s of mass at least \(\delta/2\) has
frequent \(A\)-mass at least \(\delta/2\).  Draw \(L_0\)
independent \(A\)-probes.  At each stage the spans of all past
offsets have total dimension at most \(2rL_0<d_0\).  The mass
of a probe whose whole offset list lies in those spans is
less than \(\delta/4\).  Thus each stage has conditional
probability at least \(\delta/4\) of being frequent and having
a fresh offset vector.  Select one fresh vector by a fixed
rule using only the probe lists.  Averaging fixes a probe
collection with independent selected witnesses in each
component and mode, and with successful \(B\)-mass at least
\begin{equation}
 q_*=(\delta/2)(\delta/4)^{L_0}>0.
 \label{eq:fresh-probe-mass}
\end{equation}
This lower bound is independent of selector and channel sizes.
Require
\begin{equation}
 L\,O(L_0)\,2^{L_0-h}<q_*/2.
 \label{eq:early-h-probes}
\end{equation}
By \Cref{lem:channel-thinning}, the successful \(B\)'s can
then be trimmed so that the channel
transpose at endpoint \(1\) has full rank on all these fixed
witness vectors, while leaving mass at least \(q_*/2\), after absorbing the uniform
\(o(1)\) error in a strict version of \Cref{eq:early-h-probes}.

Now estimate this event under the raw law \(\rawlaw^2\).
Condition on all of endpoint \(2\) and on the primal frames
at endpoint \(1\).  Enumerate the offset record \(m\), and
enumerate a nominal primal tuple representing \(Z_B\) at
endpoint \(1\).  There are at most \(2^{.002N+O(n)}\) such
choices.  For each choice the ambient value \(z\), all the
lists in \Cref{eq:frequent-list}, and their endpoint-\(2\)
translations are fixed.  The event defined by the adaptive
marks is contained in the union of these raw events.  This
uses no conditional density assertion about the marked law,
and does not enumerate \(2^{t_BN}\) arbitrary ambient
companion values.

Project each of the \(L_0\) lists onto the coordinate of its
chosen fresh witness.  Their total product size is at most
\(2^{.51L_0N}\).  Freeze the bounded lists of channel
coefficient values on the witnesses.  On their full-rank
event, prescribing the corresponding endpoint-\(1\) output
tuple costs at least \(.99L_0N\) bits, by the same annihilator
count as in \Cref{eq:singleton-joint-pmax}.  The upper
probability under \(\sigma\), after paying its joint density,
is therefore at most
\[
 2^{(D+.01)N+.002N+O(n)-.48L_0N+O(1)}.
\]
Choose \(M_0\) to make the \(O(n)/N\) term small.
Because \(L_0\ge100(D+10)\), this is exponentially small,
contradicting the mass \(q_*/2\).  The claim follows.

Discard the frequent events in a full-sum phase test.
For fixed records, this is a restriction on \(B\) alone once
the \(A\)-offset record is fixed.  The \(B\)-tuple point bound
is now \(2^{-(t_B+.5)N}\); the \(A\)-tuple bound is
\(2^{-(t_A-.01)N}\).  Walsh on \((t_A+t_B)N\) bits gives
\(2^{-.245N}\), before the \(2^{.004N}\) record count.
Thus every phase with \(S=\{1,2\}\) has magnitude at most
\(\delta+2^{-\Omega(N)}\).  Interchanging the units covers
the cases with \(R=\{1,2\}\).  Together with the singleton
estimates, all nontrivial means are less than \(.005\).

All retained laws have relative mass at least
\begin{equation}
 \delta_{\rm ph}=\frac{p_0\delta}{64R_0}>0.
 \label{eq:tp-phase-retained-mass}
\end{equation}
The large-cover span deletion leaves at least \(p_0/2\);
the tiny-cover restriction, transpose deletion, record restriction
and zero/nonzero split retain at least
\(p_0\delta/(32R_0)\) in total.  This proves the stated, slightly
smaller uniform bound.  The lower bounds on \(h\) depend on the
fixed marginal caps after these early restrictions.  Only after
\(h\) do we choose the thinning rate \(\rho\); the new full-sum
estimate then requires \(M_0\) large relative to \(1/\rho\).
No inverse mass depending on \(n\) or on a numerical table is paid.
\end{proof}

\begin{corollary}[A uniform phase compatibility margin]
\label{tp:phase-compatibility}
On either retained law of \Cref{thm:phase-alternative}, for each
fixed \(c\in\F\) and all sufficiently large \(n\),
\begin{equation}
 \PP\{u_{B,1}(\Delta_A)=u_{B,2}(\Delta_A)
       =u_{A,1}(\Delta_B)=u_{A,2}(\Delta_B)=c\}\ge1/32.
 \label{eq:tp-four-equal-phases}
\end{equation}
Also \(\PP\{F(A,B)+F(B,A)=0\}\ge .49\), and each single
prescribed phase bit has probability at least \(.49\).
\end{corollary}
\begin{proof}
In alternative \textup{(i)}, Fourier inversion of the four bits
gives every prescribed pattern probability at least
\((1-15\cdot .005)/16>1/32\).  The even-sum condition and each
single-bit condition have probability at least
\((1-.005)/2>.49\).
In alternative \textup{(ii)}, both full endpoint sums are zero
identically.  Among the sixteen Fourier characters, the four
with even support in each endpoint pair therefore equal one.
Their coefficients for the common pattern \((c,c,c,c)\) are all
one, for either value of \(c\).  Each of the other twelve
characters contains a singleton and has mean \(2^{-\Omega(N)}\).
Thus the probability in \Cref{eq:tp-four-equal-phases} is
\(1/4+O(2^{-\Omega(N)})\).  The remaining assertions follow
from the same zero sums and singleton means.
\end{proof}

\begin{lemma}[Injection on independent pairs]
\label{lem:injection}
Let \(\sigma_A,\sigma_B\) be two possibly different marked-unit
laws, each with pins of total dimension at most \(K\).
Suppose all four mixed endpoint marginals are at most
\(\Lambda\pi\), where \(\Lambda\) is fixed before \(h\).
Draw \(A\sim\sigma_A\) and \(B\sim\sigma_B\) independently;
the endpoints within either unit may be dependent.  For every
fixed \(\eps>0\), choosing \(h\) sufficiently large in terms
of \(K,\eps,g\) makes the probability that the pair fails a
small-table injection condition at most \(\eps\), for all
sufficiently large \(n\).  This is an unconditional estimate,
so its error can be subtracted from any independently
established compatibility mass.
\end{lemma}
\begin{proof}
It suffices to treat one component, endpoint pair and sign, and
then sum over at most \(16|\cE|\) tests.
All copies of \(A\) used below are drawn from \(\sigma_A\)
independently of \(B\sim\sigma_B\).  The argument uses only
this independence and the separate marginal bounds, and does
not require the two unit laws to coincide.  Let \(V_A\) be the image
of the individual pinned primal space \(D_{A,e}^+\cap\cB_i\),
and let \(H_B\subseteq\F^h\) be the protected coefficient space
at the opposite tested endpoint.  Both have dimension at most
\(K\).  Failure means that some \(0\ne v\in V_A\) has
\(Y_B\trans v\in H_B\).  Both the pins and \(H_B\) may depend
on their entire units.

Put \(\ell=\lfloor N/(10(K+1))\rfloor\).  The probability that
an independent new individual primal span meets a fixed space
of dimension at most \(K\ell\) is at most
\[
 \Lambda 2^{\dim\cB+K\ell-N+1}=2^{-\Omega(N)}.
\]
This uses only the unchanged primal marginal, and tests the whole
primal span, so it is uniform over the adaptive pin choices.
If a fixed \(B\) has failure probability greater than
\(\eta>0\), successive independent draws
\(A_1,\ldots,A_\ell\) can each witness failure while the spaces
\(V_{A_j}\) remain mutually independent, with conditional
probability at least \(c=\eta/2\) at every step.  The probability
of this complete test is at least \(c^\ell\).

Conversely fix any such independent spaces.  There are at most
\(2^{K\ell}\) choices of one nonzero witness in each, and the
chosen witnesses are independent.  For iid \(Y_B\) columns their
transpose images are independent uniform elements of \(\F^h\).
For a fixed \(H_B\) of dimension at most \(K\), the probability
that all these images belong to it is at most
\(2^{-(h-K)\ell}\).  There are at most
\((K+1)2^{Kh}\) possible subspaces \(H_B\), independently of
\(n\).  The original channel marginal comparison in
\Cref{lem:frame-images} therefore bounds this event under
\(\rawlaw\) by \(C_{K,h}2^{-(h-2K)\ell}\).
Here the tested list has length proportional to \(N\); we use
this absolute reference-law count, rather than the bounded-list
rank assertion of \Cref{lem:channel-thinning}.

The endpoint bound against \(\pi\) and the absolute thinning
density loss multiply this upper bound by
\(\Lambda2^{\varepsilon_{\rm ch}N}\).  Fubini and the lower
bound \(c^\ell\) give
\[
 \PP_B\{\PP_A(\mathrm{failure}\mid B)>\eta\}
 \le C_{K,h}\Lambda
 2^{-(h-2K-\log_2(1/c))\ell+\varepsilon_{\rm ch}N}.
\]
Take
\(h>2K+\log_2(1/c)+30(K+1)\), before choosing
\(\varepsilon_{\rm ch}<10^{-4}\).  The right side is
exponentially small.  The unconditional failure probability
for this test is at most \(\eta+2^{-\Omega(N)}\).
Choose \(\eta<\eps/(32|\cE|)\), sum the tests, and then take
\(n\) sufficiently large.  This proves the assertion.
\end{proof}
